\documentclass[10pt,a4paper]{amsart}
\usepackage[margin=19mm]{geometry}
\usepackage{amsmath,amssymb,mathtools}
\usepackage[scr=rsfs,cal=euler]{mathalfa}
\usepackage{xfrac}
\usepackage[unicode,hidelinks,bookmarks=false]{hyperref}
\newcommand{\ie}{i.e.\ }
\newcommand{\cf}{cf.\ }
\newcommand{\resp}{respectively\ }
\newcommand{\Subsection}{\subsection}
\providecommand{\nobreakdash}{\nobreak\hbox{-}}
\setlength{\emergencystretch}{2em}
\multlinegap0pt
\newcommand{\BR}{{\mathbb {R}}}
\usepackage{amssymb}
\usepackage{mathtools}
\newtheorem{lem}{Lemma}
\newcommand{\CC}{{\mathcal {C}}}
\def\Del{{\Delta}}
\newcommand{\Fi}{{\mathfrak {i}}}
\def\NI{\mathrm{NI}}
\def\PI{\mathrm{PI}}
\def\p{\prime}
\title{The local Gross-Prasad conjecture over $\BR$: epsilon dichotomy}
\author{Cheng Chen, Zhilin Luo}
\date{Source: arXiv:2204.01212v2 (CC BY 4.0)}
\begin{document}
\maketitle

\noindent\emph{Source and licence.} The local Gross-Prasad conjecture over $\BR$: epsilon dichotomy. Version arXiv:2204.01212v2 (CC BY 4.0), distributed by arXiv under the Creative Commons Attribution 4.0 International licence, http://creativecommons.org/licenses/by/4.0/, as recorded in the arXiv licence metadata for this exact version and shown on the abstract page https://arxiv.org/abs/2204.01212v2. Full licence notice: \texttt{licenses/arxiv-2204.01212-CC-BY-4.0.txt}.\par\smallskip

\noindent\textit{Editorial scope.} Counted unit: the article's lem lem:fiberoverkappapp (source0.tex lines 2554--2568). The article's complete original proof of it is delivered with it (source0.tex lines 2570--2611, one \texttt{proof} environment of 42 lines). No mathematics is written, repaired or re-derived: the statement and the proof are the article's own lines, in the article's own notation, and no retained line is substituted or re-flowed.  Retained uncounted as the source-local context the proof needs: lines 1349--1353.  Nothing was omitted from any retained span, so the package carries no omission record.  No retained line contains a citation, so the wrapper carries no bibliography and no bibliography entry.  No retained line contains a figure environment, an image-inclusion command or drawing code, so no image is delivered and the package declares no asset dependency.  Editorial formatting only, in addition: an AMS \texttt{amsart} preamble that restates the article's own declarations and macros used by the retained lines, namely the environments \texttt{lem} on separate counters, together with the standard packages those lines need.  The retained environments print in this document as Lemma 1 in order of appearance, whereas the article prints the same items with the numbering of its own full document.  The complete native source of the article as distributed in the arXiv e-print for this exact version is bundled unchanged as \texttt{sources/124111/source0.tex}.
\begin{equation}\label{eq:indexofWkappac}
\PI(W_{\kappa,c})=\frac{d}{2}+\theta,
\qquad
\NI(W_{\kappa,c})=\frac{d}{2}-\theta.
\end{equation}

\begin{lem}\label{lem:fiberoverkappapp}
For $\kappa^{\p\p}\in \Xi(d_V,d_W)$, the fiber of
\[
\CC(V,W)\longrightarrow \Xi(d_V,d_W)
\]
over $\kappa^{\p\p}$ is
\[
\begin{cases}
C(\kappa^{\p\p})_{\frac{\Del_W-\Fi_{W,\kappa^{\p\p}}}{2}},
& \dim W \text{ odd},\\[4pt]
C(\kappa^{\p\p})_{\frac{\Del_V-\Fi_{V,\kappa^{\p\p}}}{2}},
& \dim W \text{ even}.
\end{cases}
\]
\end{lem}

\begin{proof}
By definition, $(\kappa^{\p\p},c^{\p\p})\in \CC(V,W)$ if and only if
\[
(W_{\kappa^{\p\p},c^{\p\p}},q_{\kappa^{\p\p},c^{\p\p}})
\hookrightarrow (W,q_W)
\]
and the orthogonal complement is quasi-split. By \eqref{eq:indexofWkappac},
this is equivalent to
\[
2\sum_{i\in I^*_{\kappa^{\p\p}}}c_i^{\p\p}
=
\Del_{W_{\kappa^{\p\p},c^{\p\p}}}
=
\Del_W-\Del_{W^\perp_{\kappa^{\p\p},c^{\p\p}}}.
\]

Assume first that $\dim W$ is odd. Then the complement
$W^\perp_{\kappa^{\p\p},c^{\p\p}}$ is odd-dimensional and quasi-split, hence
\[
\Del_{W^\perp_{\kappa^{\p\p},c^{\p\p}}}
=
\Fi_{W,\kappa^{\p\p}}.
\]
Therefore
\[
\sum_{i\in I^*_{\kappa^{\p\p}}}c_i^{\p\p}
=
\frac{\Del_W-\Fi_{W,\kappa^{\p\p}}}{2}.
\]
Conversely, this equality forces the complement to have the required
quasi-split signature. This gives the first case.

If $\dim W$ is even, the same argument is applied after embedding
$W_{\kappa^{\p\p},c^{\p\p}}$ into $V$; the relevant complement is then
odd-dimensional. One obtains
\[
\sum_{i\in I^*_{\kappa^{\p\p}}}c_i^{\p\p}
=
\frac{\Del_V-\Fi_{V,\kappa^{\p\p}}}{2}.
\]
This proves the second case.
\end{proof}

\end{document}
